\section*{Hoofdstuk \ref{ch:g3:separating-mixtures} --- Mengsels scheiden}

\begin{solution}{exo:g3:separating-mixtures:1}
\omterm{def:g3:separating-mixtures:sieve}{Zeven}: de rijstkorrels zijn kleiner dan de erwten en vallen door een zeef
met de juiste gaatjes. (Met de hand sorteren lukt ook, maar langzaam.)
\end{solution}

\begin{solution}{exo:g3:separating-mixtures:2}
De aarde zakt langzaam en vormt een laag op de bodem; het water erboven
wordt helderder. Dat is \omterm{def:g3:separating-mixtures:decant}{bezinken}.
\end{solution}

\begin{solution}{exo:g3:separating-mixtures:3}
Het \omterm{def:g3:separating-mixtures:filter}{filtraat} is de vloeistof die door het filter gaat. Bij het
koffiezetten is het \omterm{def:g3:separating-mixtures:filter}{filtraat} de koffie in de kan; het koffiedik blijft in
het papier.
\end{solution}

\begin{solution}{exo:g3:separating-mixtures:4}
Ja. Het zout is opgelost, en een filter houdt niet tegen wat opgelost
is: het gaat met het water mee.
\end{solution}

\begin{solution}{exo:g3:separating-mixtures:5}
Door het water in te dampen: in de zon of zacht verwarmd gaat het water
de lucht in en blijft de suiker in het schaaltje.
\end{solution}

\begin{solution}{exo:g3:separating-mixtures:6}
Het grind blijft in de zeef. De zandkorrels zijn kleiner dan de gaatjes,
dus vallen ze erdoorheen.
\end{solution}

\begin{solution}{exo:g3:separating-mixtures:7}
(a) \omterm{def:g3:separating-mixtures:sieve}{Zeven}. (b) \omterm{def:g3:separating-mixtures:filter}{Filtreren} (of laten \omterm{def:g3:separating-mixtures:decant}{bezinken} en \omterm{def:g3:separating-mixtures:decant}{afschenken}). (c) \omterm{def:g3:separating-mixtures:evaporate}{Indampen}.
\end{solution}

\begin{solution}{exo:g3:separating-mixtures:8}
Vier stappen: roosters, bezinkbassins, zandfilters, kiemen gedood. De
bezinkbassins halen het slib weg.
\end{solution}

\begin{solution}{exo:g3:separating-mixtures:9}
\omterm{def:g3:separating-mixtures:filter}{Filtreren}; het zand van het papier halen; het \omterm{def:g3:separating-mixtures:filter}{filtraat} \omterm{def:g3:separating-mixtures:evaporate}{indampen} om het
zout te krijgen.
\end{solution}

\begin{solution}{exo:g3:separating-mixtures:10}
Nee. Het zout in zeewater is opgelost, en een filter houdt niet tegen wat
opgelost is: het \omterm{def:g3:separating-mixtures:filter}{filtraat} is nog steeds zout water.
\end{solution}

\begin{solution}{exo:g3:separating-mixtures:11}
De thee voorzichtig afgieten zodra de blaadjes bezonken zijn
(\omterm{def:g3:separating-mixtures:decant}{afschenken}), of hem door een theezeefje of een filter gieten (\omterm{def:g3:separating-mixtures:sieve}{zeven} of
\omterm{def:g3:separating-mixtures:filter}{filtreren}).
\end{solution}
